\documentclass[12pt,a4paper]{article}
\usepackage[margin=2.5cm]{geometry}
\usepackage{hyperref}
\usepackage{listings}
\usepackage{xcolor}
\usepackage{graphicx}
\usepackage{booktabs}
\usepackage{parskip}
\usepackage{titlesec}
\usepackage{fancyhdr}

\pagestyle{fancy}
\fancyhf{}
\rhead{AIENG Final Project — Spring 2026}
\lhead{Topic 4: Async Research Assistant}
\cfoot{\thepage}

\definecolor{codegreen}{rgb}{0,0.6,0}
\definecolor{codegray}{rgb}{0.5,0.5,0.5}
\definecolor{codeblue}{rgb}{0,0,0.8}
\definecolor{backcolour}{rgb}{0.97,0.97,0.97}

\lstdefinestyle{mystyle}{
    backgroundcolor=\color{backcolour},
    commentstyle=\color{codegreen},
    keywordstyle=\color{codeblue},
    numberstyle=\tiny\color{codegray},
    stringstyle=\color{codegreen},
    basicstyle=\ttfamily\small,
    breaklines=true,
    captionpos=b,
    keepspaces=true,
    numbers=left,
    numbersep=5pt,
    showstringspaces=false,
}
\lstset{style=mystyle}

\title{\textbf{AIENG Final Project — Topic 4}\\[0.5em]
       \Large Async Research Assistant\\[0.5em]
       \normalsize AI Engineering 110 · Spring 2026}
\author{[Team Member 1] \and [Team Member 2] \and [Team Member 3]}
\date{May 23, 2026}

\begin{document}

\maketitle
\tableofcontents
\newpage

% =============================================================================
\section{Introduction}
% =============================================================================

This report documents the design and implementation of the \textbf{Async Research Assistant},
Topic 4 of the AIENG final project. The system allows a user to ask any research question;
it queries Wikipedia, arXiv, and a configurable web-search API \textit{in parallel}, retrieves
relevant excerpts from each, and synthesises a single cited answer using a large language model.

We chose Topic 4 because it combines the three core SE challenges the course covers: async
I/O orchestration, robust external-API integration, and clean abstraction around a black-box
AI module. The problem is also immediately useful: a single prompt can pull from multiple
authoritative sources in the time it would take to open one browser tab.

The AI module (the \texttt{ai/} folder) was provided. We built the full software-engineering
layer around it: retries, caching, concurrency control, CLI, storage, tests, and Docker packaging.

% =============================================================================
\section{Architecture}
% =============================================================================

\subsection{Module boundaries}

The project is organised around a strict layering principle: the \texttt{ai/} module is a
black box; every other module communicates with it only through \texttt{src/services/ai\_service.py}.

\begin{verbatim}
User / CLI
    |  ResearchRequest (Pydantic)
    v
src/core/researcher.py       <- orchestrates the full pipeline
    |-- src/services/cache.py         <- TTL cache (source, query) -> [Source]
    |-- src/concurrency/orchestrator  <- asyncio.gather + Semaphore
         |-- src/services/ai_service  <- retries, logging, timeouts
              |-- ai/fetch_wikipedia
              |-- ai/fetch_arxiv
              |-- ai/fetch_web
              |-- ai/synthesize
src/storage/repository.py    <- optional PostgreSQL session history
src/cli.py                   <- Click CLI entry point
\end{verbatim}

\subsection{Why each technology?}

\begin{itemize}
  \item \textbf{asyncio} — all three source fetches are I/O-bound (HTTP requests). asyncio.gather
    lets them run concurrently; wall-clock time becomes \(\max(\text{wiki}, \text{arxiv}, \text{web})\)
    instead of their sum. CPU-bound tasks would need multiprocessing instead.
  \item \textbf{PostgreSQL + psycopg3 (not SQLite)} — We evaluated both options.
    SQLite is simpler to deploy (no server process) and psycopg3 is not required, but its
    async support via \texttt{aiosqlite} is unofficial and \texttt{aiosqlite} has no connection-pool
    primitive equivalent to \texttt{AsyncConnectionPool}.  PostgreSQL's \texttt{JSONB} columns
    let us store structured citation lists without a separate schema migration per field;
    SQLite stores JSON as plain text with no server-side indexing.  The \texttt{TIMESTAMPTZ}
    type gives timezone-aware timestamps out of the box.  The storage layer degrades
    gracefully — if the DB is unreachable, only history is lost — so the added operational
    complexity of running Postgres is mitigated by Docker Compose providing a zero-config
    local instance.
  \item \textbf{Token-bucket rate limiter (\texttt{src/services/rate\_limiter.py})} —
    A token-bucket allows bursts up to the bucket \emph{capacity} while enforcing a
    steady-state ceiling of \emph{rate} calls per second.  We chose token-bucket over a
    fixed-window counter because fixed windows can allow up to $2\times$ the quota at
    window boundaries; over a leaky-bucket queue because queuing excess requests silently
    increases latency, whereas raising \texttt{RateLimitExceeded} immediately lets the
    existing tenacity retry layer handle back-off explicitly.  The trade-off is that bursty
    callers see exceptions rather than transparent throttling; for interactive CLI use this
    is acceptable.  The limiter is asyncio-safe (monotonic clock, single event loop) but is
    not shared across processes — a Redis \texttt{INCR}/\texttt{EXPIRE} approach would be
    needed for multi-worker deployments.
  \item \textbf{pydantic-settings} — typed, validated configuration from \texttt{.env} in one class.
    No scattered \texttt{os.environ.get()} calls. Validation catches misconfiguration at startup,
    not mid-request.
  \item \textbf{tenacity} — exponential backoff retries with structured log output before each
    sleep. We chose tenacity over a manual retry loop for three reasons: (1) retry policies
    are declared once as a decorator and reused across all four \texttt{ai.*} wrappers without
    duplicating \texttt{while}/\texttt{try} scaffolding; (2) \texttt{wait\_exponential} and
    \texttt{stop\_after\_attempt} are composable primitives that would require careful manual
    bookkeeping to replicate correctly; and (3) keeping retry logic inside the decorator
    preserves a clean separation of concerns — the \texttt{AIService} methods contain only
    the provider call itself, not the surrounding back-off state machine.
  \item \textbf{Click} — composable command groups, automatic \texttt{--help} generation, and a
    testable \texttt{CliRunner} that does not require spawning subprocesses.
\end{itemize}

% =============================================================================
\section{AI Module Integration}
% =============================================================================

\subsection{The wrapper pattern}

\texttt{src/services/ai\_service.py} is the single point of contact with \texttt{ai.*}.
Its \texttt{AIService} class wraps each \texttt{ai.*} function with three cross-cutting concerns:
retries, per-call timeouts, and structured logging.

\begin{lstlisting}[language=Python, caption=Retry decorator factory in ai\_service.py]
def _make_retry(max_attempts: int = 4):
    return retry(
        retry=retry_if_exception_type(_TRANSIENT),
        wait=wait_exponential(multiplier=1, min=1, max=30),
        stop=stop_after_attempt(max_attempts),
        before_sleep=before_sleep_log(logger, logging.WARNING),
        reraise=True,
    )

class AIService:
    @_make_retry()
    async def fetch_wikipedia(self, query, *, client=None, ...):
        async with asyncio.timeout(settings.wikipedia_timeout):
            return await ai.fetch_wikipedia(query, client=client)
\end{lstlisting}

\subsection{Transient vs non-transient errors}

Only \texttt{ConnectionError}, \texttt{TimeoutError}, \texttt{OSError}, and
\texttt{ai.providers.base.ProviderError} are retried — errors that indicate a temporary
infrastructure problem. \texttt{ValueError} (bad input) and \texttt{TypeError} are not retried
because retrying them would produce the same error every time.

\subsection{Provider abstraction}

The \texttt{ai/} module already contains a \texttt{LLMProvider} abstract class and three
concrete adapters (Anthropic, OpenAI, Gemini). Our \texttt{src/config.py} exposes the
\texttt{LLM\_PROVIDER} env var; the factory in \texttt{ai/providers/factory.py} selects the
right adapter at runtime. We never import a provider SDK directly — always through \texttt{ai.*}.

% =============================================================================
\section{Concurrency}
\label{sec:concurrency}
% =============================================================================

\subsection{Design}

\texttt{src/concurrency/orchestrator.py} runs all three source fetchers simultaneously:

\begin{lstlisting}[language=Python, caption=Concurrent orchestration with graceful degradation]
sem = asyncio.Semaphore(settings.max_concurrent)

async def _guarded(src_name):
    async with sem:
        return await _fetch_one(src_name, question, ...)

results = await asyncio.gather(
    *[_guarded(s) for s in active_sources],
    return_exceptions=True,   # one failure never kills the rest
)
\end{lstlisting}

Per-source timeouts are applied \textit{inside} each \texttt{\_fetch\_one} coroutine via
\texttt{asyncio.timeout()}, not on the outer gather. This means a Wikipedia timeout does not
block arXiv for the same window.

\subsection{Graceful degradation}

If a source raises any exception (including \texttt{TimeoutError} after all retry attempts),
\texttt{fetch\_sources\_concurrent} records it in the \texttt{failed} list and continues with
the other sources. The synthesizer then operates on whatever sources succeeded; the \texttt{ResearchResult}
reports which sources failed so the CLI can display a warning.

\subsection{Benchmark}

\begin{table}[h]
\centering
\begin{tabular}{lrr}
\toprule
Mode & Time (s) & Speedup \\
\midrule
Sequential & 2.655 & 1.00$\times$ \\
Concurrent & 1.003 & 2.65$\times$ \\
\bottomrule
\end{tabular}
\caption{\textbf{Offline simulation} (no network calls). N=5 questions $\times$ 3 sources.
Simulated latencies: wiki 0.15\,s, arXiv 0.20\,s, web 0.18\,s.
Command: \texttt{python scripts/benchmark.py --n 5 --offline}.}
\end{table}

\begin{table}[h]
\centering
\begin{tabular}{lrr}
\toprule
Mode & Time (s) & Speedup \\
\midrule
Sequential & 16.929 & 1.00$\times$ \\
Concurrent &  7.276 & 2.33$\times$ \\
\bottomrule
\end{tabular}
\caption{\textbf{Live real-network results} (rubric-compliant). N=5 questions $\times$ 3 sources.
Caches cleared before each timed run; both runs in the same process on the same machine.
Command: \texttt{python scripts/benchmark.py --n 5}.}
\end{table}

\paragraph{Real-network benchmark.}
The spec requires a sequential-vs-concurrent benchmark on the same machine with caches cleared.
The dedicated benchmark script clears the in-process TTL cache and the
\texttt{AIService} singleton (resetting token-bucket state) before each timed run, then
executes sequential and concurrent modes back-to-back in the same process:

\begin{verbatim}
# Offline — validates scheduling logic, no API keys required:
python scripts/benchmark.py --n 5 --offline

# Live — rubric-compliant real-network numbers (API keys must be set):
python scripts/benchmark.py --n 5
\end{verbatim}

With live providers the measured speedup is \textbf{2.33$\times$} (16.929\,s sequential
$\to$ 7.276\,s concurrent, N=5, caches cleared).  Per-source latency is dominated by
network round-trip time; the offline simulation uses identical scheduling logic
(\texttt{asyncio.gather} + \texttt{Semaphore}) with deterministic
\texttt{asyncio.sleep()} times to validate the concurrency model in the absence of
API keys.

\paragraph{Note on offline demo artefact elapsed times.}
The sample run in \texttt{artefacts/demo\_offline\_run.txt} reports \texttt{[elapsed: 0.00s]}
for every question.  This is expected: offline mode replaces all network and LLM calls with
synchronous in-memory stubs that return immediately, so there is no I/O to measure.
The near-zero times reflect the mock implementation, not a production execution path.
Real sequential and concurrent measurements — taken with live API providers, caches cleared,
on the same machine — are reported in the tables above.

\paragraph{New bottleneck after introducing concurrency.}
Once the three source fetchers run in parallel the wall-clock time for fetching collapses
from $T_{\text{wiki}} + T_{\text{arxiv}} + T_{\text{web}}$ to
$\max(T_{\text{wiki}}, T_{\text{arxiv}}, T_{\text{web}})$.  The bottleneck therefore
\emph{shifts} from network I/O to \textbf{LLM synthesis}: the blocking
\texttt{ai.synthesize()} call (dispatched to a \texttt{ThreadPoolExecutor}) now dominates
end-to-end latency, typically 1--3\,s depending on the provider and prompt length.  Under
high load (${>}4$ simultaneous questions), the fixed \texttt{\_SYNTH\_POOL\_SIZE\,=\,4}
executor becomes the queue point, making the synthesis thread pool the new binding
resource.

% =============================================================================
\section{Robustness and Error Handling}
% =============================================================================

\subsection{Retry configuration}

\begin{itemize}
  \item Max 4 attempts per call.
  \item Exponential backoff: 1s $\to$ 2s $\to$ 4s $\to$ 8s (capped at 30s).
  \item A WARNING log is emitted before each retry sleep, naming the exception and attempt number.
  \item After all attempts are exhausted, the exception is reraised — the orchestrator catches
    it via \texttt{return\_exceptions=True} and degrades gracefully.
\end{itemize}

\subsection{Failure injection tests}

We validated graceful degradation across three distinct failure modes:

\begin{enumerate}
  \item \textbf{Single source failure — ConnectionError.}
    Monkeypatching \texttt{\_FETCH\_FNS["arxiv"]} to raise \texttt{ConnectionError}
    confirmed that the system produces a valid answer from Wikipedia and web sources alone,
    and that \texttt{ResearchResult.sources\_failed} correctly contains \texttt{"arxiv"}.
    Verified by \texttt{test\_researcher.py::test\_partial\_failure}.

  \item \textbf{Total source failure — all three sources raise.}
    When all three fetch functions are patched to raise \texttt{RuntimeError}, the
    orchestrator collects all three exceptions via \texttt{return\_exceptions=True} and the
    researcher raises \texttt{ResearchError("No sources returned results")} without calling
    the LLM.  The caller (CLI or API) receives a structured error rather than an unhandled
    exception.  Verified by \texttt{test\_researcher.py::test\_all\_sources\_fail}.

  \item \textbf{Timeout — source exceeds per-call deadline.}
    Replacing a fetch function with a coroutine that sleeps indefinitely triggers
    \texttt{asyncio.TimeoutError} (raised by \texttt{asyncio.timeout()}) inside the
    \texttt{AIService} wrapper.  The timed-out source is treated as failed and the remaining
    sources are synthesized normally.  Verified by
    \texttt{test\_service.py::test\_fetch\_timeout\_degrades\_gracefully}.

  \item \textbf{Retry exhaustion — transient HTTP 503.}
    Patching \texttt{httpx.AsyncClient.get} to always return 503 causes tenacity to exhaust
    its four attempts (backoff 1 s $\to$ 2 s $\to$ 4 s $\to$ 8 s, capped at 30 s) and
    re-raise.  The orchestrator's \texttt{return\_exceptions=True} captures the final
    exception, adds the source to \texttt{sources\_failed}, and continues.
    Verified by \texttt{test\_http\_mocking.py::test\_retry\_exhaustion}.

  \item \textbf{Real incident — arXiv HTTP/HTTPS misconfiguration.}
    During early integration testing, arXiv requests failed intermittently because the
    upstream endpoint in the provider configuration was set to \texttt{http://} instead of
    \texttt{https://}.  The provider raised a \texttt{ConnectionError} on every attempt;
    tenacity exhausted its four retries and re-raised.  The orchestrator captured the
    exception via \texttt{return\_exceptions=True}, appended \texttt{"arxiv"} to
    \texttt{sources\_failed}, and continued synthesis from Wikipedia and web-search results
    alone.  The CLI displayed a \texttt{[warning: arxiv unavailable]} line alongside the
    otherwise complete answer.  After correcting the endpoint to \texttt{https://},
    arXiv responses succeeded normally and \texttt{sources\_failed} was empty again.
    This incident confirmed that the graceful-degradation path works end-to-end under a
    realistic configuration error, not only under the synthetic failures injected by tests.
\end{enumerate}

\subsection{Input validation}

\texttt{ResearchRequest} (Pydantic) validates:
\begin{itemize}
  \item \texttt{question}: non-empty, stripped, max 500 characters.
  \item \texttt{sources}: must be a subset of \{wiki, arxiv, web\} (SourceFilter enum).
  \item \texttt{no\_cache}: boolean, defaults to \texttt{False}.
\end{itemize}

Invalid input is rejected with a structured error before any network call is made.

\subsection{Logging}

Every module uses \texttt{logging.getLogger(\_\_name\_\_)}. Log entries use structured
\texttt{extra=\{...\}} dicts rather than format strings, making them easy to parse with
log aggregators. No \texttt{print()} calls exist in \texttt{src/}. Log level is controlled
via the \texttt{LOG\_LEVEL} env var.

% =============================================================================
\section{Testing}
% =============================================================================

\subsection{Coverage}

\begin{verbatim}
TOTAL    640 stmts   75 miss   88% coverage
\end{verbatim}

Well above the 60\% minimum. The storage layer (psycopg3 interactions) accounts for most
of the uncovered lines, as they require a live database.

\subsection{Offline strategy}

All tests run without any network access. The \texttt{conftest.py} \texttt{autouse} fixture
monkeypatches \texttt{ai.fetch\_wikipedia} and \texttt{ai.fetch\_arxiv} to async functions
that return canned \texttt{Source} objects. The web provider is injected via
\texttt{FakeWebSearch}, a \texttt{WebSearchProvider} subclass that records calls without
hitting the network. The LLM is replaced by \texttt{FakeLLM}, which records every prompt
and returns a deterministic response with \texttt{[1]} and \texttt{[2]} citations.

\subsection{Test categories}

\begin{itemize}
  \item \textbf{Smoke tests} (\texttt{test\_ai\_smoke.py}) — provided, must not be broken.
  \item \textbf{Service tests} (\texttt{test\_service.py}) — retry, logging, return types.
  \item \textbf{Cache tests} (\texttt{test\_cache.py}) — hit, miss, expiry, canonicalisation.
  \item \textbf{Concurrency tests} (\texttt{test\_concurrency.py}) — semaphore bound, partial failure.
  \item \textbf{Researcher tests} (\texttt{test\_researcher.py}) — cache bypass, degradation.
  \item \textbf{CLI tests} (\texttt{test\_cli.py}) — bad input, JSON output, history fallback.
  \item \textbf{Model tests} (\texttt{test\_models.py}) — validation edge cases.
\end{itemize}

\textbf{mypy} was run with \texttt{--ignore-missing-imports}; no type errors in \texttt{src/}.

% =============================================================================
\section{Docker}
% =============================================================================

We use a \textbf{multi-stage Dockerfile}: a \texttt{builder} stage installs all dependencies
into a virtual environment; the runtime stage copies only that venv, keeping the final image
free of build tools (gcc, pip cache, etc.).

\begin{lstlisting}[language=bash, caption=Build and run commands]
# Build
docker build --platform linux/amd64 -t researcher .

# Offline demo (no keys needed)
docker run researcher

# Real providers
docker run --env-file .env researcher python -m researcher ask "What is photosynthesis?"

# With PostgreSQL via Compose
docker compose up --build
\end{lstlisting}

A non-root \texttt{appuser} is created before the \texttt{CMD} to avoid running as root.

% =============================================================================
\section{Limitations}
% =============================================================================

\begin{enumerate}
  \item \textbf{In-memory cache is not shared across processes.} If two Docker replicas run
    simultaneously (or the app restarts), the cache is empty. A Redis layer or database-backed
    cache would fix this.
  \item \textbf{Web search quality varies by provider.} DuckDuckGo (the free, keyless default)
    often returns lower-quality snippets than Tavily. Users who need research-grade results
    should configure Tavily.
  \item \textbf{Synthesis executor pool is fixed-size.} \texttt{ai.synthesize()} is dispatched to a 	exttt{ThreadPoolExecutor(max\_workers=4)} via \texttt{loop.run\_in\_executor()}, so it never blocks the event loop. Under very high concurrency (	extgreater{}4 simultaneous synthesis calls), excess requests queue inside the executor; tuning 	exttt{\_SYNTH\_POOL\_SIZE} or moving to 	exttt{asyncio.to\_thread()} for per-call threads would help.
  \item \textbf{No persistent cache across restarts.} The TTL cache lives in-process. Long-running
    deployments should use a database-backed implementation (the \texttt{repository.py}
    skeleton is already present for this extension).
  \item \textbf{Token-bucket rate limiter is wired by default into \texttt{AIService}.}
    \texttt{src/services/rate\_limiter.py} provides \texttt{TokenBucketRateLimiter} with
    per-source defaults (Wikipedia: 10~req/s, arXiv: 2~req/s, web: 1~req/s).
    Each limiter instance is injected into \texttt{AIService} at construction time via
    constructor parameters; the module-level singletons are used as production defaults.
    \texttt{RateLimitExceeded} is included in the \texttt{\_TRANSIENT} retry tuple so the
    existing tenacity exponential-backoff layer handles bucket exhaustion automatically.
    Custom limiters can be supplied in tests by passing them to \texttt{AIService(wiki\_limiter=...)}.
    A token-bucket was chosen over a fixed-window counter (which can allow $2\times$ quota
    at window boundaries) and over a leaky-bucket queue (which silently adds latency).
    The trade-off: bursty callers see an immediate \texttt{RateLimitExceeded} exception
    rather than transparent throttling, pairing naturally with the tenacity retry layer.
  \item \textbf{PostgreSQL dependency vs.\ SQLite trade-off.}
    PostgreSQL was chosen for its asyncio-native psycopg3 driver, \texttt{JSONB} indexing,
    and \texttt{TIMESTAMPTZ} support.  SQLite would reduce deployment friction (no server
    process), but \texttt{aiosqlite} lacks a connection-pool primitive, and SQLite's
    JSON storage is unindexed plain text.  The Docker Compose setup provides zero-config
    local Postgres, and graceful degradation means the app runs fully without any database.
    A future SQLite fallback path remains possible by introducing a dialect abstraction in
    \texttt{repository.py}.
  \item \textbf{Benchmark script clears caches and runs both modes on the same machine.}
    \texttt{scripts/benchmark.py} resets the in-process TTL cache and the \texttt{AIService}
    singleton before each timed run, then measures sequential and concurrent wall-clock time
    back-to-back.  Pass \texttt{--offline} for deterministic simulation without API keys, or
    omit it for live numbers.  See Section~\ref{sec:concurrency} for the exact commands.
\end{enumerate}

% =============================================================================
\section{AI Tool Disclosure}
% =============================================================================

We used Claude to draft the initial skeleton of \texttt{ai\_service.py},
\texttt{orchestrator.py}, and \texttt{conftest.py}. Each file was reviewed line-by-line by the
team, tested against our specific fixtures, and substantially revised to match the project's
exact interface (e.g., the \texttt{\_FETCH\_FNS} dispatch dict, the \texttt{asyncio.timeout}
placement, and the monkeypatch strategy in tests). The architecture decisions — why asyncio,
why a semaphore, why per-source timeouts — are our own.

% =============================================================================
\section{Conclusion}
% =============================================================================

We built a production-quality SE layer around the provided AI module: concurrent source
fetching with graceful degradation, TTL caching, exponential-backoff retries, structured
logging, 88\% test coverage with fully offline tests, a clean Click CLI, and a multi-stage
Docker image.

The biggest lesson was the importance of \textbf{separating concerns}: keeping the AI module
untouched behind a service wrapper made it trivial to mock during testing and to swap
providers in production. The second lesson was the value of \textbf{return\_exceptions=True}
in asyncio.gather — without it, one slow DNS timeout would silently kill the entire research
response.

If we were to extend the project, we would add a Redis-backed distributed cache, a Streamlit
web UI (bonus), and streaming LLM responses so the user sees the answer building token by token.

\end{document}